% 来源：OEBPS/Text/Section0001_split_001.xhtml#sigil_toc_id_12

\BookChapter{4}
\section[第4章\q 掌控金钱：进行突破]{第4章\\ 掌控金钱：进行突破}

\begin{BookQuote}
感恩是高贵灵魂的标志。

\hfill ——伊索
\end{BookQuote}

钱是一种工具，能让我们把梦想变成现实，没有足够的钱或者说特别缺钱，你的生活会很悲惨。但是当你口袋里有了足够的钱，是不是所有的事情就会自动变好了呢？我想我们都知道答案是什么。

我们是什么人，钱无法改变。钱能做的只是放大我们的本性，如果你本性吝啬又自私，有了钱你会更加吝啬，更加自私。如果你感恩又博爱，有了钱你会更多地感恩他人，更多地给予他人。

我们稍停片刻，回顾一下2008年的金融危机。房地产市场和股票市场暴跌，数万亿美元的市值都灰飞烟灭，化为乌有。短短几个月，数百万人失业。你经历了什么呢？你遭受的打击重吗？这对你的家庭影响大吗？你的朋友怎么样了？我们每个人的反应不同，有的人恐惧，有的人愤怒，有的人辞职，有的人决心改变。所有这些反应都与钱无关，而是关于我们自己的。这些反应透露出我们的金钱观——钱对于我们来说到底意味着什么，我们赋予钱什么样的力量。要么我们让钱控制自己，要么我们自己控制钱。

\subsection{要钱还是要生活}

在2008年金融危机期间，有很多人、很多事让人震惊，我个人认为影响最大的案例来自一个名叫阿道夫·默克尔的绅士。2007年，他在全世界富豪排行榜上的排名是第94位，也是德国最富有的人，个人拥有财富120亿美元。他拥有欧洲最大的制药企业，后来又把自己的商业帝国拓展到制造业和建筑业。他很骄傲自己能拥有这么大的成就。他也是一个相当厉害的投机者。

2008年，阿道夫·默克尔决定在股票市场上大赌一把。他非常肯定，大众汽车公司的股票一定会下跌，所以他决定做空这只股票。但是半路出了一个意外：保时捷公司收购了大众汽车，大众汽车的股价暴涨，并没有像他预料的那样下跌。几乎一夜之间，默克尔输掉了7.5亿美元。

雪上加霜的是，阿道夫·默克尔迫切地需要一些现金来偿付巨额债务，但是2008年，银行不会贷款给任何人：不贷给你，不贷给我，也不贷给亿万富翁，甚至也不贷给其他银行。

默克尔是怎么应对的呢？寻找新的融资？减少开支？亏本卖掉一些公司的股票救急？这些事他都没有做。他发现，自己总共亏掉了30亿美元，已经不再是德国最富有的人了，他让自己的家人失望了，于是写了一份自杀声明，然后跑向正在疾驰的列车，一头撞了过去。

没错，他自杀了。

具有悲剧性讽刺意味的是，他的家人发现，默克尔自杀后，才过了几天，他申请的贷款就被批下来了，他的企业保住了。

阿道夫·默克尔是因为钱才死的吗？或者说，他选择死亡是因为钱对于他来说意义很大吗？对于阿道夫·默克尔来说，钱就是身份。没有钱，就没有身份，没有地位。失去了德国首富这个地位，他完全无法承受，他觉得自己简直太失败了——尽管亏了30亿美元后，他口袋里还有90亿美元！

你可能会想：“哎，太可惜了。”但是对于我们这些旁观者来说，这样判断阿道夫·默克尔这个人，可能太武断了。但是想想我们自己，我们是不是也经常把自己的身份，自己未来的前途，在某种程度上和钱紧紧联系在一起？我们不愿意承认这一点，但其实是这样的。

\subsection{想要破产死掉的亿万富翁}

另一方面，也有像查克·费尼那样的人，他是一个美籍爱尔兰人，出生在新泽西州的伊丽莎白市，是一个白手起家的亿万富翁。你在全世界任何一个地方，在机场里穿行时，发现自己受到吸引走进一家商店，里面都是装着酒和香水的闪闪发光的瓶子，还有其他的免税奢侈品，你是否有这样的体验呢？免税店。对的。免税店就是查克·费尼的创意。他1960年从一无所有开始创业，白手起家，创立了一个免税店销售的企业帝国，身家75亿美元。

一度，《福布斯》世界富豪排行榜把他列为世界首富，但是查克·费尼非常低调，你可能从来没有听说过他当上世界首富这个新闻。他一生大部分时间中，个人名下的资产既没有汽车，也没有房子。他坐飞机都是坐经济舱，戴的手表是塑料手表。和阿道夫·默克尔一样，他的银行账户里的钱也缩水了很多——现在他已经80多岁了，自己名下的财产只有100多万美元。但是他和默克尔有非常大的不同，他没有像默克尔那样想要紧紧地抓住每一分钱，查克·费尼把自己所有的钱都捐了出去。

就是这个家伙，最近30年来，把自己的人生使命确定为利用钱这个工具，来改变各地方人的生活。他的慈善活动遍及全世界，在北爱尔兰帮助创造和平，在南非帮助人们与艾滋病斗争，在芝加哥帮助教育儿童。

查克·费尼这个人最让人意想不到的事情是，他做所有这些慈善活动都是匿名的。费尼不想得到任何荣誉。事实上，直到最近才有消息透露出来，他就是所有这些不可思议的计划背后捐款支持的那个人。而且他还在继续这样匿名做着好事。查克·费尼说他的目标是要做慈善活动做到花光最后一分钱为止，做到他签的最后一张支票银行拒绝兑现为止。

很明显，对于阿道夫·默克尔和查克·费尼来说，钱的意义完全不同。那么金钱到底对于你来说的真正意义是什么呢？是你用钱还是钱来用你？就像我一开始所说的那样：如果你不能掌控钱，在某种程度上，钱就会掌控你。

\subsection{终极目标：回馈社会}

对于我本人来说，当我还是小孩子的时候，钱总是看不见摸不着的。一提钱，我们就头痛，因为我们家的钱很少，从来不够花。我记得很清楚，我不得不去敲邻居家的门，借点食物给我的兄弟姐妹和我自己吃。

后来我11岁那年，在感恩节那天（类似于中国人的春节），发生了一件事，由此永远改变了我的人生。那一天跟往年的感恩节一样，家里什么吃的都没有，我的父母两人还在打架。突然我听见有人在敲门，我把门打开一条缝，看见一个人站在台阶上，手里提着个篮子，里面装满了食物，足够做上一顿感恩节大餐。我简直不敢相信自己的眼睛。

我的父亲总是说，没有人会白白地送别人东西。但是，突然有个我根本不认识的人，什么回报也不要，过来关心我们，给我们送温暖。这件事让我想到：“这是不是我听说过的那种来自陌生人的关怀？”我决定，如果有陌生人能够关心我和我的家庭，给我们送温暖，我将来也会关心那些陌生人。“我将来应该怎么做呢？”那一天我向自己承诺，我要找到一种办法，用某种方式，在某一天，去回馈社会，把爱心传递下去。就这样，等我长到17岁的时候，我晚上兼职当看门的保安，攒了些钱，在感恩节那天出去给两个陌生的家庭送温暖，送给他们一大堆食物。这是我一生中最感动的经历。看到那些面孔从绝望变成欢乐，这让我感到自己的精神境界提升了。真的，这件事对于他们来说是一份大礼，对于我来说同样也是一份大礼。我没有告诉任何人我做了这件好事，第二年我给四个家庭送去了食物。再到后来，我给八个家庭送去了食物。我做这些不是为了得到别人的夸奖、赞美，不是像小学生那样为了评上“优秀学生”而做好人好事。但是，我给八个家庭送去了食物之后，我发现做这样的好事，我一个人确实能力有限。我想：“兄弟，我可以找些人来帮我一起做。”于是我召集了一些朋友，他们也加入了我们，逢年过节一起来给困难的家庭送食物、送温暖。我们捐助的人数越来越多。现在我的慈善基金会举办的“国际食物篮”慈善行活动，每年给36个国家200万人送食物、送温暖。要不是我11岁的那个神奇的感恩节有陌生人给我们家送食物，我怎么会知道给予他人会如此快乐？谁知道呢？有些人会说，这就是运气，这就是命运，或者是那个很古老的说法，这就是福气。我看到了上帝之手，我称之为“恩典”。

下面是我领悟到的东西：我懂得了给予的快乐，这和金钱毫无关系，金钱只是一种工具，用来满足我们的需求，而不仅是我们的财务需求。我们的人生之路大部分受到信念的指引，而这些信念是在人生路上随着时间推移而逐步形成的；我们的人生之路大部分也受到故事的指引，这个故事是由我们创造的，人生为了什么，我们本来应该成为什么样的人，我们本来应该做什么，给予什么。最终，是什么让我们感到快乐和满足？每个人都有不同的快乐，有些人觉得快乐就是愉悦他人，有些人觉得得到权力和控制别人才会快乐，有些人把自己的快乐定义为得到亿万财富，有些人觉得找到快乐和有意义的人生道路就是和上帝更加接近，放弃所有的物质东西，还有一些人觉得快乐的最终意义就是自由。

不管你追求的是什么样的情绪，不管你追逐的是什么样的工具（创立一家企业，养一家子人，环游世界），不管你想要的天堂是什么，我发现这些都是你的一种企图，你费尽心机想要满足人类六大需求中的一种或几种。

激发我们行动的就是这六大基本需求。所有的人类行为，其背后都是由这六大需求驱动的，全世界都是这样。无论人们做的是疯狂的坏事，还是伟大的好事，行为背后的推动力都是这六大需求，但是这六大需求在我们心中的价值地位不同，满足需求的顺序也不同，这就决定了我们人生的方向不同。

为什么人类六大基本需求非常重要，我们必须深入理解它们？因为你要创造财富，就必须知道你真正追求的是什么——你积累这些财富到底是为了什么。你积累财富是为了让自己觉得确定而且安全吗？你追求财富是为了让自己觉得很特别、与众不同吗？你追求的是一种奉献的感觉——你想为别人做点事情，用你过去从来不能用的方式吗？还是以上这些动机你都有呢？

如果你把确定性看作你生命中最重要的需求，和你把爱情作为第一大需求相比，确定性会让你的行动方向有很大的不同，在人际关系上、在商业上、在投资理财上也会有很大不同。你如果深入挖掘自己真正追求的东西，就会发现那其实根本不是钱。你真正追求的是你认为钱能够给你带来的东西。那是一系列感觉的组合。在这些感觉背后是你的基本需求。

\subsection{人类第一大需求：确定性/舒适感}

人类的第一大需求是确定性。我们需要控制感，知道接下来会发生什么事，这会使我们感觉安全。这是基本舒适性的需求，是避免痛苦和压力的需求，同时也是创造欢乐的需求。这样的需求正常吗？确定性需求是我们的一种生存本能。它会影响我们在生活中愿意冒多大的风险——在我们的工作中，在我们的投资上，在我们的人际关系上。对确定性的需求越强，你实际愿意承受的风险就越小，或者说你的情绪愿意忍受的风险就越小。顺便说一下，这就是你真正的“风险承受力”的来源。

但是你如果任何时候都能完全确定，又会怎么样？你知道什么事情会发生，什么时候会发生，会怎样发生，你在别人开口说话之前就知道他们要说什么，这会让你有什么样的感觉？刚开始你会感觉自己很神奇、很特别，但是最终你会感觉怎么样？你会感觉太没有意思了。

\subsection{人类第二大基本需求：不确定性/变化}

上帝啊，您的智慧无限大，才会给我们人类第二大基本需求，它就是不确定性。我们需要变化，我们需要惊奇。

我来问你一个问题：你喜欢惊奇吗？

如果你回答“是”，你就是在戏弄你自己！你想得到的惊奇，你称之为惊喜。你不想得到的惊奇，你称之为问题。但是你还是需要那些让你惊奇的问题来给你的生活增添一些活力。你不能锻炼你的肌肉或是你的性格，除非你有东西在后面压着你的肌肉或性格。没有压迫就没有反抗，也就没有成长。

\subsection{人类第三大基本需求：重要性}

第三大需求是重要性。正是这个基本需求驱动阿道夫·默克尔努力打拼，成为德国首富，也是这个需求驱动他失去首富地位后自杀了。我们都需要感觉自己很重要，很特别，与众不同，别人很需要我们。但是我们如何得到重要性？有很多得到别人尊重的方式，赚到亿万财富能得到别人尊重，获得很高的学位——硕士、博士，能让人觉得你有文化、与众不同。你可以成为一个“大V”（微博上有大群粉丝的用户），有很多粉丝追随，让自己感到很受尊重。你可以参加《单身汉》，或者参加下一季《贵妇的真实生活》，一夜成为大明星。有些人得到别人重视的方式是在身上到处刺文身穿孔，或者在私密的地方刺文身穿孔。你要得到别人的重视还有另一种方式，就是你的问题比别人更多更大：“你是不是觉得我的丈夫就是一袋大垃圾？求求你把我丈夫这袋垃圾拿走一天吧。”当然你得到别人的尊重也可以通过在精神层面上高人一等（假装也行）。不幸的是，要得到别人重视的最快方式，既不需要用钱，也不需要读书、受教育，那就是用暴力。如果有人拿枪指着你的头，那么就在那一刻，他马上变成你生命中最重视的人，对不对？

花很多钱，能让你自己感觉自己很重要，但花很少的钱你也能做到。我们都知道，人们不停地吹嘘他们买的东西多么便宜。有人觉得自己很特别，因为他们用牛粪烧火取暖，或者只借助太阳能来取暖。有些非常富有的人之所以得到别人的重视，是因为他们隐藏自己的巨额财富。山姆·沃顿，沃尔玛超市的创始人，有很长一段时间都是美国首富。他自己开着车在阿肯色州本顿维尔跑来跑去，那辆车还是一辆旧皮卡，这表明他根本不需要一辆宾利豪车，不过他确实有一架私人飞机随时可用。

重要性，也是一个赚钱机器。这正是我的亲密好友史蒂夫·韦恩创造财富的秘密。就是这家伙打造了拉斯维加斯这个赌城。他知道人们愿意买单，只要他们相信这是最好的东西——让他们觉得自己很特别、与众不同、很重要。史蒂夫·韦恩在他自己开的赌场和酒店里提供你能想象到的独家、奢侈豪华的体验。他的赌场和酒店雄伟华丽，无与伦比。他搞了一个夜总会叫作“XS”（还有什么型号呢），那里是拉斯维加斯最受欢迎的地方。即使是在平时工作日的晚上，“XS”门口也排着长队。你一进去，就有特权可以购买一瓶普通的香槟，一瓶700美元，如果你想显摆一下，显示你是个大玩家，你可以花10000美元买一瓶特制的“小野鸡尾酒”，它用罕见的白兰地和新鲜的橙汁调制而成，杯子上还挂着一条白金项链。嘿，你看，服务员送到你桌子上时，它还喷着焰火，照亮了整个大厅。你就是要让所有人都知道你很重要，很受重视（这只是你自己的想象而已）。

\subsection{人类第四大基本需求：爱和联系}

人类第四大基本需求是爱和联系。爱对于生命来说就像氧气对人体一样重要。爱是我们最想要得到的，也是我们最需要得到的。当我们完全投入去爱的时候，我感觉到自己活着，但是当我们失去爱的时候，痛苦会非常大，很多人只能靠和别人的联系才能活下来，尽管这些联系只是爱的碎屑。你可以得到那种与人联系的感觉，或者爱的感觉，通过亲密关系，通过真诚的友谊，通过祈祷，通过在大自然中行走。如果这些都没有效果，你可以去养一条狗。

前面四种人类基本需求，我称之为人格的需求。我们都可以找到方法来满足这四种需求：工作更加努力，解决大问题，创造出故事来把这些需求合理化。最后两个基本需求是精神需求，它们更加少见——并不是每个人都能得到这些精神需求的满足。一旦满足了这些精神需求，我们会真正地感到满足。

\subsection{人类第五大基本需求：成长}

人类第五大基本需求是成长。你如果现在停止成长，会是什么样子呢？你就死了。停止成长就意味着死亡。如果两个人的关系不再成长，如果一个企业不再成长，如果你这个人不再成长，那么不管你现在在银行里有多少钱，不管你有多少朋友，不管有多少人爱你，你都不会体验到真正的满足。我们成长的理由，我相信，就是我们有一些很有价值的东西可以给予别人。

\subsection{人类第六大基本需求：奉献}

人类第六大基本需求是奉献。尽管它听起来好像是老生常谈，但人生的秘密就在于给予。人生的目的不是为了我，而是为了我们。好好想想，当你得到一个好消息，令人激动的消息，你会做的第一件事是什么？你会打电话给你爱的人，分享这个好消息。分享强化了你能体验到的每一件事。

生命就在于创造意义。生命的意义并不在于你得到了什么，而在于你给予了什么。最终，你得到的东西从来不会给你带来长久的快乐，但是你成为什么样的人，你奉献了什么，会给你带来长久的快乐。

因为本书谈的是钱，请你好好想想，钱能够满足你的六大基本需求吗？钱能满足你的确定性需求吗？你觉得肯定能。钱能满足你的变化性需求吗？检查一下。很明显，钱能让我们感觉自己很重要，很受重视。但是钱能够满足你对人际关系和爱的需求吗？正如甲壳虫乐队那首经典的歌曲所唱的那样，钱不能买来爱，但是钱能买来狗！而且不幸的是，钱能给你一种错误的联系感，因为钱能够吸引很多人和你联系，尽管并不都是那种非常让人满足的人际关系。钱能带来成长吗？钱能推动企业的成长，也能推动学习的进步。你的钱越多，你能在财务上奉献社会的就越多。

但是我真正相信的是，如果你把重要性看得高于其他所有需求，那么你想通过赚很多钱来让大家都重视你，最后会竹篮打水一场空，除非你把自己赚的钱都捐赠给社会，以此来获得重视。如果你想从钱上获得重视，那么你要付出的代价会非常高。你追求的是更大的财富数字，但是你想因此得到更大的满足感，那是根本不可能的。

生命中最终的重要性并非来自某些外在的东西，而是来自某些内在的东西。生命中最终的重要性来自我们自己对自己的尊重，而自尊可不是我们可以从别人那里得到的。别人可以告诉你，你很漂亮，你很聪明，你很明智，你最优秀；他们也可以告诉你，你是地球上最可怕的家伙。但是真正重要的是你觉得自己怎么样。你是不是相信你内心深处、灵魂深处还在继续成长，继续推动你努力前进，去做更多，去给予更多，不满足于你现在的水平，甚至不满足于你认为自己能够达到的最高水平。

没有东西会比成长和给予更加重要。尽管钱是一个特别厉害的工具，可以满足我们的六大基本需求，但是钱只是工具的一种。当你在追求钱的时候，不要忘记你为什么要追求钱。你其实是想满足某些情绪上和心理上的需求。在这些情绪的背后，是那些想让你的人生与众不同而必须满足的基本需求。

宇航员在月球上行走时，想象一下他们走过的人生旅程。他们从还是孩子时就梦想着有一天能够飞到外太空。直到那一天，巴兹·奥尔德林和尼尔·阿姆斯特朗站在月球上，回望我们居住的地球的奇观。我们大家只是在照片上看过地球的全貌，而他们两个人是人类历史上第一批亲眼看到地球全貌的人——这是多么重要的伟大成就啊。

但是后来他们怎么样了？他们参加了空中抛撒彩色纸带的大游行，受到了总统接见。他们成了美国家喻户晓的英雄人物。后来又怎么样了？假使你已经在月球上行走过了，现在只有39岁，那么以后你要做些什么？如果你研究过宇航员的历史，读过他们写的传记，你就会知道，很多宇航员都变得特别沮丧。为什么？因为他们能找到的唯一冒险方式就是太空旅行或者飞向月球。他们忘记了如何从日常生活中一个简单的微笑里找到冒险的快乐。

我不想给你更多说教，我只想用这短暂的时间来告诉你，现在是时候去掌控你的钱，而不是等着让钱来掌控你。要找到最快的方式来满足你的六大基本需求，感觉到你和别人的联系，感觉到你的人生多么重要，感觉到确定感和变化，把你自己置于一种给予别人的状态。每天找到一种方式，让自己感激别人更多，期盼自己得到更少，地球上最富有的人就是那个能够感激别人更多的人。

我第一次采访约翰·邓普顿爵士，是在我33岁的时候。他是身家几十亿美元的超级富豪，但是他起步时一无所有，他是白手起家的，而他能赚到如此巨大的财富，都是在历史上最糟糕的年代，别人都感到非常恐惧的时候：第二次世界大战后的日本；20世纪80年代末到90年代初，南美洲部分国家受到了巨大通货膨胀的冲击。当其他人非常恐惧的时候，邓普顿勇敢地站出来，大量投资买入。我问他：“财富的秘密是什么？”他说：“托尼，你知道的，你非常清楚地知道，你可以把这个财富的秘密教给每一个人，它就是感激。”你满怀感激的时候，就没有恐惧。你满怀感激的时候，就没有愤怒。在我认识的人里面，约翰·邓普顿爵士是最快乐、最满足的一个人。尽管他在2008年去世了，但这些年来他的人生经历继续激励着所有的人。

如果你想非常富有，那么现在开始吧。今天你会对什么充满感激？今天你会对谁充满感激？你会不会对自己一路经历的那些问题和痛苦充满感激？如果你有了新的信念，你相信自己生命中发生的每一件事，都是为了服务你、帮助你，你又会怎么想？如果你从内心深处相信，人生种种经历并非偶然发生在你身上，而是专门为你发生的，你又会怎么想？你人生之路的每一步都是在帮助自己变得更加强大，让自己能做得更好，让自己能享受到更多，让自己能给予更多。如果你从这一点开始，钱就不再是你快乐的来源，也不再是你痛苦的来源。赚钱将只是一个被你逐步熟练掌握的有趣旅程，财富就变成了一个伟大的工具，让你成就那些生命中最重要的事情。

不过，既然钱是我们生命中如此重要的一部分，让我们先回到赚钱的轨道上。尽管我们这一章说的都是发自肺腑的真心话，但是你在投资理财之路上遇到的很多人，并不是从自我成长和奉献他人的角度来做投资交易的！你要进入一个竞争激烈的金融世界，里面的很多人和机构都在找机会欺负你，占你的便宜，因为你是一个缺少经验、不懂投资的菜鸟。所以，我想让你提前做好充分准备。在我们开始讨论把钱投资到哪里和寻找投资机会之前，我要告诉你注意避免哪些投资误区。

有一个原因可以解释为什么大多数投资者长期来看并没有赚到钱。我想给你必要的知识，既能保护你的投资安全，又能让你抓住机会让自己的投资增长最大化，从而能够让你实现真正的财务自由，速度比你想象的还要快。你值得拥有心灵的平静，很快就能真的拥有。翻到下一页……
